

\chapter*{Oświadczenie o wykorzystaniu narzędzi sztucznej inteligencji}
\addcontentsline{toc}{chapter}{Oświadczenie o wykorzystaniu narzędzi sztucznej inteligencji}

Podczas przygotowania niniejszej pracy korzystano z narzędzi generatywnej sztucznej
inteligencji. Zakres ich wykorzystania obejmował:

\begin{itemize}
  \item generowanie i refaktoryzację kodu źródłowego implementującego metody
        odszumiania oraz zunifikowany potok treningowy,
  \item wykrywanie błędów i optymalizację istniejącego kodu,
  \item przekład opisów algorytmicznych i pseudokodu z publikacji źródłowych
        na implementacje w języku Python,
  \item wsparcie redakcyjne przy tekście pracy oraz korektę językową,
  \item objaśnianie zagadnień teoretycznych z zakresu cyfrowego przetwarzania sygnałów.
\end{itemize}

Wykorzystane narzędzia: \textit{UZUPEŁNIĆ --- nazwy, wersje i okres użycia}.

Narzędzia sztucznej inteligencji nie były wykorzystywane do formułowania wniosków,
interpretacji wyników pomiarów, doboru pozycji bibliograficznych ani do generowania
danych pomiarowych. Wszystkie wyniki liczbowe pochodzą z obliczeń przeprowadzonych na
danych zarejestrowanych samodzielnie oraz na publicznie dostępnych bazach sygnałów.
Każda pozycja bibliograficzna została zweryfikowana w~źródle oryginalnym.

Kod powstały przy wsparciu narzędzi sztucznej inteligencji został poddany weryfikacji
obejmującej kontrolę zgodności implementacji z~opisem w~publikacji źródłowej, testy
poprawności wykonania oraz porównanie wyników z~wartościami wyznaczonymi analitycznie
dla przypadków o~znanym rozwiązaniu. Szczegółowy opis procedury weryfikacji zawiera
rozdział~\ref{ch:implementacja}.

Autor ponosi pełną odpowiedzialność za treść pracy, poprawność implementacji,
rzetelność przedstawionych wyników oraz zgodność cytowań ze źródłami.
